\documentclass[10pt]{article}
\usepackage[letterpaper,margin=0.7in,top=0.6in,bottom=0.6in]{geometry}
\usepackage[T1]{fontenc}
\usepackage{microtype,array,booktabs,enumitem,amsmath,amssymb,xcolor}
\usepackage[numbers,sort&compress]{natbib}
\usepackage[colorlinks,allcolors=blue!50!black]{hyperref}
\usepackage[compact]{titlesec}
\titleformat{\section}{\large\bfseries}{\thesection}{0.5em}{}
\titlespacing*{\section}{0pt}{3pt}{1pt}
\setlist{nosep,leftmargin=1.2em}
\setlength{\parindent}{0pt}\setlength{\parskip}{2pt}
\setlength{\bibsep}{0pt}
\renewcommand{\bibfont}{\footnotesize}
\newcommand{\acc}{I already have a GAIVI account}
\newcommand{\Gt}{G_t}

\begin{document}
\begin{center}
{\Large\bfseries Soft Value-Based Decoding for Guided Molecular Generation\\[-1pt] on a Frozen Graph-Diffusion Model}\\[3pt]
Abel Le \quad$\cdot$\quad CIS 6930 Generative AI, Fall 2026 \quad$\cdot$\quad \textbf{Enrollment: graduate (6930)} \quad$\cdot$\quad Path A (own topic)
\end{center}
\vspace{-4pt}

\section{Question}
Reward-guided generation usually fine-tunes the generator (e.g., RL fine-tuning~\cite{black2024ddpo}), which is costly. Soft Value-Based Decoding (SVDD)~\cite{li2024svdd} instead steers a frozen diffusion model at inference time: at every reverse step it draws $M$ candidates, scores each with a value function, and keeps one. Its molecule experiments use a continuous model (GDSS) and defer DiGress to future work~\cite[App.~C]{li2024svdd}. Whether SVDD works on a \emph{discrete graph} diffusion model, and whether it beats drawing the same number of candidates as independent samples and keeping the best, is therefore open.
\textbf{This project asks whether SVDD with a learned value function, applied to a frozen DiGress model trained on QM9, raises the mean synthetic-accessibility and QED reward of generated molecules more than best-of-$N$ sampling with $N{=}M$ independent samples per slot, without lowering validity, and how that advantage scales with the number of candidates $M$.}

\section{Primary Paper and Background}
\textbf{Primary:} SVDD~\cite{li2024svdd} shows that a soft value function $v_t(x_t)=\alpha\log\mathbb{E}[e^{r(x_0)/\alpha}\mid x_t]$ lets a pretrained diffusion model sample from a reward-tilted distribution with no fine-tuning, on images, DNA/RNA and molecules (default $M{=}10$, argmax selection).
\textbf{Supporting:} DiGress~\cite{vignac2023digress} is a discrete graph diffusion model whose denoiser predicts clean atom and bond types (QM9: 99.0\% valid, 96.2\% unique, about 1\,h to reach 99\% validity). DDPO~\cite{black2024ddpo} is the fine-tuning route that SVDD avoids. Gao et al.~\cite{gao2023overopt} study best-of-$n$ as a simple, strong reward-optimization baseline. TreeDiff~\cite{zhao2026treediff} benchmarks SVDD and best-of-$N$ on a different graph model (GraphDiT) under matched compute.
\textbf{Gap:} the SVDD paper's molecule results use GDSS on ZINC-250k and defer DiGress to future work. I found no study of SVDD on DiGress/QM9 that counts network passes, compares against best-of-$N$ over several seeds, and diagnoses the learned value function. I measure that. This is new work started this semester, not prior research.

\section{Method}
\textbf{Pipeline (all inference-time; DiGress weights frozen).} (1)~\emph{Proposal:} the official QM9 checkpoint link is dead (404; repository issue~\#109), so I train DiGress once on QM9 with the repository's unmodified QM9 config (cross-entropy on node/edge types, $T{=}500$) and then freeze it (one training run that only provides the frozen base model). A thin wrapper exposes $p_\theta(G_{t-1}\mid\Gt)=\sum_{G_0}q(G_{t-1}\mid\Gt,G_0)\,p_\theta(G_0\mid\Gt)$ and returns $M$ i.i.d.\ candidates per molecule per step; one network pass gives $p_\theta(G_0\mid\Gt)$ for all $M$ draws. (2)~\emph{Value function} $\hat v_\phi(\Gt,t)$, trained by MSE regression onto $r(G_0)$ using 10k unguided rollouts (the SVDD-MC estimator); two architectures: an MLP on pooled atom/bond-type histograms plus a timestep embedding, and a 3-layer message-passing GNN on the (soft) graph. (3)~\emph{Selection:} keep $\arg\max_m \hat v_\phi(G_{t-1}^{(m)})$ (or sample $\propto e^{\hat v/\alpha}$). A training-free variant, SVDD-PM, scores a candidate by $r(\hat G_0)$ of its decoded denoiser prediction (one extra pass per candidate). (4)~\emph{Rewards:} SA reward $(10-\mathrm{SA})/9$ from the SA score~\cite{ertl2009sa} and QED~\cite{bickerton2012qed} via RDKit, invalid $\to 0$. Value functions are trained per reward. \textbf{Stretch:} a docking-proxy reward, dropped if not started by Nov 20 (see Risks).

\section{Data}
QM9~\cite{ramakrishnan2014qm9}: 133{,}885 molecules (up to 9 heavy atoms; the MoleculeNet CSV that DiGress downloads), public, CC0 on figshare, no personal data. DiGress trains on its own seeded 100k/val/test split of QM9 (repository code); afterwards I only \emph{generate}. All of QM9 is the novelty reference. Three seeded (42), disjoint 1k index splits (committed) only sanity-check reward code and value-function predictions at $t{=}0$ on real molecules, never training or tuning. Value-function data are unguided rollouts from the frozen DiGress (separate training/validation chains and random streams from evaluation). Each evaluation run generates $N{=}2{,}000$ molecules.

\section{Evaluation Plan}
\textbf{Metrics:} mean and top-10\% reward, validity (RDKit-sanitizable and connected; stricter than DiGress's largest-fragment convention, which I also report), uniqueness, novelty vs.\ QM9, diversity (mean pairwise ECFP4 Tanimoto distance), network passes and wall-clock per molecule. \textbf{Seeds:} 0, 1, 2 for every run, plus 3, 4 for the headline point ($M{=}10$); report mean$\pm$std. \textbf{Baselines:} (B1) unguided DiGress ($M{=}1$); (B2) \emph{best-of-$N$}, $N{=}M$: draw $M$ unguided molecules per slot and keep the highest-reward one (true reward, a strong baseline). It matches SVDD in candidates but spends up to $M\times$ more network passes, so I also report reward against network passes and wall-clock, including best-of-$N$ at SVDD's wall-clock; (B3) SVDD-PM without a learned value. \textbf{Controlled comparisons:} (A1)~$M\in\{1,5,10,20\}$ for SA and QED; (A2)~value estimator $\in\{$PM, MLP, GNN$\}$ at $M{=}10$. Not run: DDPO-style fine-tuning (outside the budget) and DiGress's regressor guidance (needs a retrained regressor; a different question).
\textbf{Falsifiable outcome:} the hypothesis is wrong if, at $M{=}10$, SVDD-MLP's mean-reward gain over B2 is not larger than twice the across-seed standard deviation (the larger of the two methods'), \emph{or} if SVDD reaches its gain only with validity more than 5 points below B1. A null or negative result will be reported and analyzed.

\section{Compute Plan}\small
\textbf{Where:} GAIVI (class partition, Slurm; scripts in repo), Colab/local as backup; \acc{}; first job in week 1. \textbf{Estimate} per run $=N p/512\times0.1\,\mathrm{s}\times1.5/3600$ h, with $N{=}2000$ molecules, $p$ network passes per molecule ($T{=}500$ for unguided and SVDD-MLP, $(M{+}1)T$ for SVDD-PM, $MT$ for best-of-$M$), batches of 512 graphs, an assumed 0.1 s per batch-pass ($\le$9-node graphs; \emph{measured in week 1}) and 1.5$\times$ overhead for featurizing and scoring; one $p{=}T$ run costs 0.08 h.
B1 (5 seeds): 0.4 h. B2: one pool of 20 draws per slot per seed yields best-of-5/10/20 for both rewards: $3\times1.6+2\times0.8=6.5$ h. A1 (24 runs $\times0.08$ h, plus per (reward, seed) 11k rollout chains 0.45 h and one fit 0.5 h, assumed): 7.6 h. A2: SVDD-PM at $M{=}10$ and a random-selection control 5.9 h; GNN arm 6.2 h. Seeds 3--4 for SVDD-MLP at $M{=}10$: 4.1 h. Subtotal 30.7 h; with 25\% margin 38 h; plus DiGress training on QM9, budget $\le$10 GPU-h (paper: about 1 h to 99\% validity): \textbf{$\approx$48 GPU-hours}.
\textbf{Single-GPU fallback ($<$2\,h):} $N{=}1000$, $M\in\{1,10\}$, one seed, MLP value, SA reward, on a short-trained DiGress checkpoint: two runs 0.04 h each, 1.2k rollout chains 0.05 h, one fit 0.5 h. Training checkpoints are kept and sampling is saved after every chunk of 512 molecules (a preempted run resumes to identical molecules); every run logs config, seed, git commit, GPU and metrics.

\normalsize\section{Milestones (approval Oct 12 $\to$ presentation, targeting Slot A, Fri Nov 13)}
\begin{tabular}{@{}p{1.15in}p{5.7in}@{}}
Oct 12--18 & GAIVI environment (Python 3.9, torch 2.0.1); pin commits; first GPU job; \textbf{launch DiGress QM9 training}; measure s/epoch, s/batch-pass. Sign up for Slot A on Oct 13.\\
Oct 19--25 & Finish training and reproduce DiGress validity ($\ge$97\%, paper 99.0\%); finish the \texttt{propose()} wrapper with exactness tests; B1 over 3 seeds. \textbf{Gate Oct 25.}\\
Oct 26--Nov 1 & Collect rollouts; train MLP value fn and check it at $t{=}0$; SVDD-PM end to end; B2; complete the $<$2\,h fallback run. \textbf{Gate Fri Oct 30:} go/no-go for Slot A (else ask for Slot C).\\
Nov 2--8 & \textbf{Main comparison at $M{=}10$} (starts Oct 31), SA and QED, 3 seeds: SVDD-MLP vs.\ B2, B1, SVDD-PM; start the $M$ sweep (A1). Optional Nov 6 office-hours check-in. Results frozen Nov 8.\\
Nov 9--13 & Slides from frozen results (due Nov 12), rehearse code walk-through; \textbf{present Nov 13} (pipeline, main result, baselines; ablations shown as planned/partial).\\
Nov 14--22 & Finish A1 (all $M$, both rewards); A2 incl.\ GNN value fn (\textbf{go/no-go Nov 20}); seeds 3--4 at $M{=}10$.\\
Nov 23--29 & Failure analysis (where the value fn mis-ranks), figures, seed variance checks; results frozen Nov 29.\\
Nov 30--Dec 11 & Report and clean-clone check; \textbf{report and code due Dec 11}.\\
\end{tabular}

\section{Risk and Fallback}\small
(i)~If my DiGress QM9 model is below 97\% validity (DiGress's metric) by Oct 25, I train longer or try the repository's \texttt{fixed\_bug} branch (total extra budget $\le$10 GPU-h); if it stays below 90\%, I freeze it as is and report B1's lower validity (the 5-point rule is relative to B1). (ii)~If the GNN value function does not beat the MLP on validation MSE by Nov 20, I drop it and keep PM vs.\ MLP as ablation A2; the docking proxy is dropped if not started by Nov 20. (iii)~If by Oct 18 the measured s/batch-pass is more than 3$\times$ my estimate, I drop seeds 3--4, then cut $N$ to 1000 and the $M$ grid to $\{1,10,20\}$ so the total stays near 50 GPU-h. (iv)~If by Nov 8 SVDD does not beat B2, the deliverable becomes a negative-result analysis of the value function's ranking error. (v)~If the Oct 30 gate fails, I ask to present in Slot C (Dec 4), same order of work.

\normalsize\section{Contribution Table (full version in \texttt{CONTRIBUTIONS.md})}
\small
\begin{tabular}{@{}>{\raggedright\arraybackslash}p{1.15in}>{\raggedright\arraybackslash}p{2.75in}>{\raggedright\arraybackslash}p{3.0in}@{}}
\toprule
Component & Origin & What I changed\\\midrule
DiGress model & \texttt{cvignac/DiGress} @ \texttt{bdfa10c} (MIT) & Code and default QM9 config unmodified (no official checkpoint available); I train it once, then freeze; added a candidate-proposal wrapper (\texttt{vgmg/models/digress.py})\\
SVDD algorithm & Li et al.~\cite{li2024svdd}; \texttt{hyanan16/SVDD-molecule} @ \texttt{455e4d4} (GDSS/ZINC-250k, no license) read for reference only & Own implementation from the paper (\texttt{vgmg/rl/svdd.py}); ported to discrete graph diffusion\\
Value functions & Mine & MLP, GNN, and PM estimators; training loop\\
Rewards, metrics & RDKit QED; SA score from RDKit Contrib & Own wrappers and rescaling; validity/novelty/diversity code is mine\\
Best-of-$N$, ablation harness & Mine & Compute-matched baseline; sweep scripts and logging\\
\bottomrule
\end{tabular}
\normalsize

\section{AI Tool Declaration}
I used Claude Code to scaffold the repository, draft unit tests, and edit this proposal. The research question, experimental design, baselines, and ablation choices are mine; I will write and understand the DiGress wrapper and SVDD code, and all reported numbers will come from my own runs.

{\small
\bibliographystyle{plainnat}
\bibliography{refs}}
\end{document}
